% ============================================================
% base/templates/style.tex
%
% 这份是「新建 deck 时复制到目标目录」的样式文件。
%
% 它自带**内置基线的全部内容**（页面、颜色、字体、字号、页标题、间距、
% 全部组件），所以即使不和 skill 目录在一起也能编译出正确的排版。
% 内容分别对应：
%   base/templates/format.tex      页面 / 颜色 / 字体 / 字号 / 页标题 / 间距
%   base/templates/components.tex  全部组件
%
% ★ 本文件是**唯一事实来源**。format.tex / components.tex 是留在 skill 里的
%   同内容副本，供"要往基线里加东西"时对照用，**不复制进 deck** ——
%   否则它们会与 <deck>/slide/.slide-forge/templates/ 下的同名文件撞名，
%   两个同名文件都得改，迟早改分叉。
%   改了本文件里的基线内容，记得同步那两个副本。
%
% ★ 末尾会额外加载本 deck 的本地层（`.slide-forge/`），见第八节。
%   它不在时，本文件就是完整的基线。
%
% ============================================================
% 加载链（两层的编译侧）
% ------------------------------------------------------------
%   1. 本文件 = 内置基线
%   2. <deck>/slide/.slide-forge/templates/{format,components}.tex
%      本 deck 自己的一份（锻造的 + 从仓库取用的）。存在才加载。
%
% 后加载的覆盖先加载的。见第八节的详细说明。
%
% ★ 全局 evolved/ **不在**这条链上：它是仓库，不是编译层。
%   要用里面的模板就把那份**复制**进 .slide-forge/templates/。
%
% ★ 用 \IfFileExists 而不是 \input：没有本地内容的 deck 照常编译，
%   不会因为缺文件报错。
% ★ 第 2 步是相对路径，cwd 必须是 <deck>/slide/
%   （TeX 的 \input 按当前工作目录解析，不是按本文件的位置）。
% ============================================================

% ============================================================
% 一、页面与导航
% ============================================================
\usetheme{default}
\usefonttheme{professionalfonts}
\usepackage{tabularx}   % 表格用它让表宽自适应
\usepackage{graphicx}   % 图片页
\setbeamertemplate{navigation symbols}{}
\setbeamertemplate{headline}{}
\setbeamertemplate{footline}{}
\setbeamersize{text margin left=0.90cm,text margin right=0.90cm}

% ============================================================
% 二、颜色
% ------------------------------------------------------------
% 正文一律黑色，页标题用强调色；全篇不使用斜体。
% 要强调某个词时用 \textbf 加粗，一句话里不要超过一处。
% ============================================================
\definecolor{SlideInk}{HTML}{14161A}       % 正文：近黑（唯一正文色）
\definecolor{SlideAccent}{HTML}{26308F}    % 页标题：低饱和深靛蓝
\definecolor{SlideBG}{HTML}{FBFAF7}        % 暖白背景

\setbeamercolor{background canvas}{bg=SlideBG}
\setbeamercolor{normal text}{fg=SlideInk,bg=SlideBG}
\setbeamercolor{frametitle}{fg=SlideAccent,bg=SlideBG}
\setbeamercolor{itemize item}{fg=SlideInk}
\setbeamercolor{itemize subitem}{fg=SlideInk}
\setbeamercolor{structure}{fg=SlideAccent}

% ============================================================
% 三、字体
% ------------------------------------------------------------
% 中文：微软雅黑　西文与数字：Segoe UI
% 微软雅黑没有真斜体，所以全篇不使用 \emph / \textit。
% ============================================================
\setCJKsansfont{Microsoft YaHei}
\setsansfont{Segoe UI}
\renewcommand{\familydefault}{\sfdefault}

% ============================================================
% 四、字号与行距：全篇只有两级
% ------------------------------------------------------------
% 页标题档 / 正文档。正文里的一切（含姓名、单位、要点）都用
% \SlideTextFont，不存在第三级字号。
% 每一级的第二个数字就是行距，不再另外设 \linespread。
% ============================================================
\newcommand{\SlideTitleFont}{\fontsize{20}{26}\selectfont}  % 页标题档
\newcommand{\SlideTextFont}{\fontsize{12}{17}\selectfont}   % 正文档（唯一）

% 封面档：只用于封面与致谢页的主标题。
% 全篇唯一的例外 —— 封面上的姓名如果和正文一样大就不成其为封面，
% 但它只出现在"封面 / 致谢"这两页，正文页一律不用。
\newcommand{\SlideCoverFont}{\fontsize{26}{32}\selectfont}

\setbeamerfont{frametitle}{size=\SlideTitleFont,series=\bfseries}
\setbeamerfont{normal text}{size=\SlideTextFont,series=\mdseries}
\setbeamerfont{itemize/enumerate body}{size=\SlideTextFont,series=\mdseries}

% ============================================================
% 五、页标题
% ------------------------------------------------------------
% 左上角，无分隔线、无页脚、无装饰。整份 deck 靠它对齐。
% ============================================================
\setbeamertemplate{frametitle}{%
  \nointerlineskip%
  \vspace*{0.50cm}%
  \begin{beamercolorbox}[wd=\paperwidth,leftskip=0.90cm,rightskip=0.90cm,ht=0.72cm,dp=0.06cm]{frametitle}%
    \usebeamerfont{frametitle}\insertframetitle%
  \end{beamercolorbox}%
}

% 页标题与正文之间的统一间距。
% 必须用 \vspace*（带星号）：普通 \vspace 出现在段落起点时会被
% TeX 在换页处丢弃，结果就是"设了间距却没生效"。
\newlength{\SlideBodyTopSkip}
\setlength{\SlideBodyTopSkip}{0.40cm}

% ============================================================
% 六、行距与垂直节奏
% ============================================================
\setlength{\parindent}{0pt}
\setlength{\parskip}{0pt}

% 列表不用符号，改用悬挂缩进：某条折行时第二行与首行文字左对齐，
% 条目边界依然清楚。圆点在 12pt 下要么发灰、要么与基线打架。
\setlength{\leftmargini}{0pt}
\setbeamertemplate{itemize item}{}
\setbeamertemplate{itemize subitem}{}
\setbeamertemplate{itemize/enumerate body begin}{\SlideTextFont}

% ============================================================
% 七、组件
% ------------------------------------------------------------
% 页面文件只允许使用这一节的命令 + 纯文本。
% 出现 \fontsize / \vspace / \setbeamercolor 就是设计债，
% 应该在这里加组件，而不是在页面里写样式。
%
% --- 在 .slide-forge/ 里定义组件的专用命令（必读）---------------------
% `.slide-forge/templates/components.tex` 在本文件之后加载，所以它能覆盖
% 基线里的同名组件。但 \newcommand 遇到已存在的命令会**直接报错**：
%
%     ! LaTeX Error: Command \DeckX already defined.
%
% 而 \renewcommand 又要求命令已经存在 —— 单用哪一个都不行：
% 新组件用前者、覆盖用后者，可你事先不知道是哪种情况。
% 所以在本层定义组件**一律用 \DeckDefine**，它两种情况都吃：
%
%     \DeckDefine{\DeckTimeline}[2]{...}    ✅ 新组件、覆盖基线都行
%     \newcommand{\DeckTimeline}[2]{...}    ❌ 覆盖时报 already defined
%
% \DeckDefine 等价于 \providecommand 后再 \renewcommand：命令不存在就
% 新建，已存在就替换。参数写法与 \newcommand 完全一致。
% ============================================================
\newcommand{\DeckDefine}[1]{%
  \providecommand{#1}{}%
  \renewcommand{#1}%
}

% --- 有标题页骨架：标题由 main.tex 给出，正文来自页面文件 ---
\newcommand{\DeckSlide}[2]{%
  \begin{frame}[t]{#1}%
    \vspace*{\SlideBodyTopSkip}%
    \input{#2}%
  \end{frame}%
}

% --- 无标题页骨架：封面与致谢页用 ---
% #1 内容文件，#2 垂直对齐（c 居中 / t 顶对齐）
% 做法是在这一页把 frametitle 模板清空（模板是局部的，不影响其他页），
% 内容用 \vfill 上下撑开，从而垂直居中。
\newcommand{\DeckBare}[2]{%
  \begin{frame}[#2]%
    \setbeamertemplate{frametitle}{}%
    \input{#1}%
  \end{frame}%
}

% --- 正文行：全页只有这一种行 ---
% 姓名、单位、每一条要点，全部用同一个命令，因此字号、行距、
% 字重、颜色必然完全一致 —— 从代码层面就不可能出现不一致。
% 不要为"姓名"或"要点"另建命令：两个命令意味着两套样式。
\newcommand{\DeckLine}[1]{%
  \par
  \hangindent=1.1em\hangafter=1
  \noindent #1\par
  \vspace{0.30cm}%
}

% --- 分组留白 ---
% 唯一用来分组的手段。不用分隔线、不用小标题、不用字号差异。
\newcommand{\DeckGap}{%
  \par\vspace*{0.40cm}%
}

% --- 带时间标签的行 ---
% #1 时间标签，#2 内容
% 不用空格去凑对齐（标签字数不等，手打空格永远对不齐）；
% 用固定宽度的盒子把标签撑到统一宽度，内容自然左对齐。
\newlength{\DeckTagWidth}
\setlength{\DeckTagWidth}{5.6em}   % 约容纳 5 个全角字符
\newcommand{\DeckTagged}[2]{%
  \DeckLine{\makebox[\DeckTagWidth][l]{#1}#2}%
}

% --- 条目：一条两行（标题行 + 说明行）---
% #1 标题行（可内嵌 \DeckLink），#2 说明行（角色 / 出处）
% 两行之间 0.12cm，让它们读作一条完整记录；条目之间 0.34cm。
\newcommand{\DeckCite}[2]{%
  \par
  \hangindent=1.1em\hangafter=1
  \noindent #1\par
  \vspace{0.12cm}%
  \noindent #2\par
  \vspace{0.34cm}%
}

% --- 超链接 ---
% beamer 已加载 hyperref，直接用 \href。
% 颜色用强调色，不加下划线边框（hyperref 的默认方框很难看）。
\newcommand{\DeckLink}[2]{%
  \href{#1}{\color{SlideAccent}#2}%
}

% --- 表格行距（可调旋钮）---
% 默认值 = 正文行距。某页需要更紧或更松时在该页开头改：
%     \setlength{\DeckTableLeading}{14pt}
% 字号不在这里调：字号只有两档，全篇一致，不允许逐页改动。
\newlength{\DeckTableLeading}
\setlength{\DeckTableLeading}{\baselineskip}

% --- 数值表：「名称占满 + 数值右对齐」---
% #1 表体：名称 & 数值 \\
% 表宽自动等于 \linewidth。不画竖线、不画外框。
\newcommand{\DeckCourseTable}[1]{%
  \par\noindent
  {\SlideTextFont
   \setlength{\baselineskip}{\DeckTableLeading}%
   \setlength{\tabcolsep}{0pt}%
   \renewcommand{\arraystretch}{1.0}%
   \noindent
   \begin{tabularx}{\linewidth}{@{}>{\raggedright\arraybackslash}X@{\hspace{0.16cm}}r@{}}
     #1
   \end{tabularx}}%
  \par
}

% --- 列表式两列表：「窄标签左 + 长文字占满左对齐」---
% 年份 + 描述、时间线这类内容用这个，
% 否则描述会被右对齐得零碎，且长条目会溢出列宽。
\newcommand{\DeckListTable}[1]{%
  \par\noindent
  {\SlideTextFont
   \setlength{\baselineskip}{\DeckTableLeading}%
   \setlength{\tabcolsep}{0pt}%
   \renewcommand{\arraystretch}{1.0}%
   \noindent
   \begin{tabularx}{\linewidth}{@{}l@{\hspace{0.55cm}}>{\raggedright\arraybackslash}X@{}}
     #1
   \end{tabularx}}%
  \par
}

% 表头行：与正文同字号，只靠加粗区分
\newcommand{\DeckCourseHead}[2]{%
  \textbf{#1} & \textbf{#2} \\[0.04cm]
}
\newcommand{\DeckListHead}[2]{%
  \textbf{#1} & \textbf{#2} \\[0.04cm]
}
% 列头与第一行数据之间的细线。
% 用 \noalign 单独发 \hrule，避免 \hline\noalign 给表格带来的异常深度
% （实测那一手会让 5 行表的深度变成 77pt）。
\newcommand{\DeckCourseRule}{%
  \noalign{\vspace{0.09cm}\hrule height 0.4pt\vspace{0.13cm}}%
}

% --- 两图并排 ---
% #1 左图路径，#2 右图路径，#3 左图高度，#4 右图宽度
%
% 不用 beamer 的 columns：它的栏间距会与给定宽度叠加，很难算准。
% 改成两图放在同一行、用 \hfill 撑开，宽度完全可控。
%
% \includegraphics 的 width/height 是"缩放后"尺寸：只给 height 时，
% 宽度按原图比例推算，所以两张图若比例不同、都想等高，
% 就必须给其中一张显式 width（调用方算好传入）。
%
% 关键约束：两图宽度之和必须明显小于正文宽（14.25cm），留 0.3cm 以上。
% 只差 0.03cm 都会让第二张图被挤到下一行，并报 overfull vbox。
\newcommand{\DeckFigures}[4]{%
  \par\noindent
  \includegraphics[height=#3]{#1}%
  \hfill
  \includegraphics[width=#4]{#2}%
  \par
}

% ============================================================
% 八、项目本地层（.slide-forge/）—— 可选，存在才加载
% ------------------------------------------------------------
% 加载顺序 = 优先级（后加载的覆盖先加载的）：
%
%   1. 内置基线       本文件（上面一到七节）
%   2. 项目本地       .slide-forge/templates/
%                     ← 这个 deck 自己的一份。skill 自动沉淀的新模板在这里，
%                       从全局 evolved/ 取用的模板也复制到这里。
%
% ★ 只有两层。全局 evolved/ **不在这里加载** —— 它是仓库，不是编译层。
%   原因：若把 evolved/ 垫在每份 deck 底下，任何一次吸收都会**静默改掉
%   所有已有 deck 的排版**（实测过：零错误零警告，但版式变了）。
%   已交付的 deck 必须冻结，所以每份 deck 只认自己 .slide-forge/ 里的那份。
%
%   要用 evolved/ 里的模板时：**复制**到 .slide-forge/templates/ 再用，
%   复制是逐字复制，**不要顺手改内容** —— 一改就与仓库那份分叉。
%
% ★ 两行都是 \IfFileExists：文件不存在就跳过。全新的 deck 照常编译，
%   不会因为缺文件报错。
% ★ 相对路径，cwd 必须是 <项目>/slide/（TeX 的 \input 按工作目录解析）。
% ============================================================
\IfFileExists{.slide-forge/templates/format.tex}%
  {% ============================================================
% base/templates/format.tex   内置格式配置（基线 · 仓库里的源）
%
% 这一层只放"长什么样"：页面、颜色、字体、字号、页标题、间距。
% 不含任何组件（组件在 components.tex）。
%
% 【这份文件不参与 deck 编译】
%   新建 deck 时只复制 `style.tex`、`main.tex`、`pages/`。
%   本文件与 components.tex 留在 skill 里，是给 deck**抄**的源：
%   deck 的 style.tex 自带等价内容，并另加载本 deck 的 .slide-forge/ 覆盖。
%   —— 不复制它们进 deck，是为了避免与
%   <deck>/slide/.slide-forge/templates/format.tex 同名共存，
%   那种"两个同名文件都得改"的坑迟早会把内容改分叉。
%
% 【三层架构里各自的位置】
%   base/templates/style.tex          内置基线 → 复制进 deck，唯一的事实来源
%   base/templates/format.tex|components.tex   本文件，仓库里的基线源
%   <deck>/slide/.slide-forge/templates/       本 deck 的本地层（自动生成）
%   evolved/templates/format.tex|components.tex 吸收来的仓库（不参与编译）
%
% ★ 不要直接改本文件来调格式。要改请写进对应 deck 的
%   .slide-forge/templates/format.tex，满意后再走「吸收」进 evolved/。
% ============================================================

% ============================================================
% 一、页面与导航
% ============================================================
\usetheme{default}
\usefonttheme{professionalfonts}
\usepackage{tabularx}   % 表格用它让表宽自适应
\usepackage{graphicx}   % 图片页
\setbeamertemplate{navigation symbols}{}
\setbeamertemplate{headline}{}
\setbeamertemplate{footline}{}
\setbeamersize{text margin left=0.90cm,text margin right=0.90cm}

% ============================================================
% 二、颜色
% ------------------------------------------------------------
% 正文一律黑色，页标题用强调色；全篇不使用斜体
% （中文没有真斜体字形，合成倾斜会发虚）。
% 要强调某个词时用 \textbf 加粗，一句话里不要超过一处。
% ============================================================
\definecolor{SlideInk}{HTML}{14161A}       % 正文：近黑（唯一正文色）
\definecolor{SlideAccent}{HTML}{26308F}    % 页标题：低饱和深靛蓝
\definecolor{SlideBG}{HTML}{FBFAF7}        % 暖白背景

\setbeamercolor{background canvas}{bg=SlideBG}
\setbeamercolor{normal text}{fg=SlideInk,bg=SlideBG}
\setbeamercolor{frametitle}{fg=SlideAccent,bg=SlideBG}
\setbeamercolor{itemize item}{fg=SlideInk}
\setbeamercolor{itemize subitem}{fg=SlideInk}
\setbeamercolor{structure}{fg=SlideAccent}

% ============================================================
% 三、字体
% ------------------------------------------------------------
% 中文：微软雅黑　西文与数字：Segoe UI
% 微软雅黑没有真斜体，所以全篇不使用 \emph / \textit。
% 可用字重只有 Regular 与 Bold 两档 —— 正好对应"正文 / 强调"。
% ============================================================
\setCJKsansfont{Microsoft YaHei}
\setsansfont{Segoe UI}
\renewcommand{\familydefault}{\sfdefault}

% ============================================================
% 四、字号与行距：全篇只有两级
% ------------------------------------------------------------
% 页标题档 / 正文档。正文里的一切（含姓名、单位、要点）都用
% \SlideTextFont，不存在第三级字号。
% 每一级的第二个数字就是行距，不再另外设 \linespread。
% ============================================================
\newcommand{\SlideTitleFont}{\fontsize{20}{26}\selectfont}  % 页标题档
\newcommand{\SlideTextFont}{\fontsize{12}{17}\selectfont}   % 正文档（唯一）

% 封面档：只用于封面与致谢页的主标题。
% 全篇唯一的例外 —— 封面上的姓名如果和正文一样大就不成其为封面，
% 但它只出现在"封面 / 致谢"这两页，正文页一律不用。
\newcommand{\SlideCoverFont}{\fontsize{26}{32}\selectfont}

\setbeamerfont{frametitle}{size=\SlideTitleFont,series=\bfseries}
\setbeamerfont{normal text}{size=\SlideTextFont,series=\mdseries}
\setbeamerfont{itemize/enumerate body}{size=\SlideTextFont,series=\mdseries}

% ============================================================
% 五、页标题
% ------------------------------------------------------------
% 左上角，无分隔线、无页脚、无装饰。整份 deck 靠它对齐。
% ============================================================
\setbeamertemplate{frametitle}{%
  \nointerlineskip%
  \vspace*{0.50cm}%
  \begin{beamercolorbox}[wd=\paperwidth,leftskip=0.90cm,rightskip=0.90cm,ht=0.72cm,dp=0.06cm]{frametitle}%
    \usebeamerfont{frametitle}\insertframetitle%
  \end{beamercolorbox}%
}

% 页标题与正文之间的统一间距：整页呼吸感的起点。
% 必须用 \vspace*（带星号）：普通 \vspace 出现在段落起点时会被
% TeX 在换页处丢弃，结果就是"设了间距却没生效"。
\newlength{\SlideBodyTopSkip}
\setlength{\SlideBodyTopSkip}{0.40cm}

% ============================================================
% 六、行距与垂直节奏
% ============================================================
\setlength{\parindent}{0pt}
\setlength{\parskip}{0pt}

% 列表不用符号，改用悬挂缩进：某条折行时第二行与首行文字左对齐，
% 条目边界依然清楚。圆点在 12pt 下要么发灰、要么与基线打架。
\setlength{\leftmargini}{0pt}
\setbeamertemplate{itemize item}{}
\setbeamertemplate{itemize subitem}{}
\setbeamertemplate{itemize/enumerate body begin}{\SlideTextFont}
}{}
\IfFileExists{.slide-forge/templates/components.tex}%
  {% ============================================================
% base/templates/components.tex   内置组件（基线 · 仓库里的源）
%
% 页面文件只允许使用本文件里的命令 + 纯文本。
% 出现 \fontsize / \vspace / \setbeamercolor 就是设计债，
% 应该加组件，而不是在页面里写样式。
%
% 【这份文件不参与 deck 编译】
%   新建 deck 时只复制 `style.tex`、`main.tex`、`pages/`。
%   本文件与 format.tex 留在 skill 里，是给 deck**抄**的源：
%   deck 的 style.tex 自带等价内容，并另加载本 deck 的 .slide-forge/。
%
% 【要新增组件写到哪】
%   <deck>/slide/.slide-forge/templates/components.tex   ← 自动生成，不用问用户
%   evolved/templates/components.tex                     ← 吸收，用户点头后才做
%   本文件：只有用户明确同意"升入基线"时才动。
% ============================================================

% ============================================================
% 页面骨架
% ============================================================

% --- 有标题页：标题由 main.tex 给出，正文来自页面文件 ---
\newcommand{\DeckSlide}[2]{%
  \begin{frame}[t]{#1}%
    \vspace*{\SlideBodyTopSkip}%
    \input{#2}%
  \end{frame}%
}

% --- 无标题页：封面与致谢页用 ---
% #1 内容文件，#2 垂直对齐（c 居中 / t 顶对齐）
%
% 这两页不要页标题：封面本来就该是整页一块，致谢页只有两个字，
% 再顶一个页标题会显得多余、也会和正文页混淆。
% 做法是在这一页把 frametitle 模板清空（模板是局部的，不影响其他页），
% 内容用 \vfill 上下撑开，从而垂直居中。
\newcommand{\DeckBare}[2]{%
  \begin{frame}[#2]%
    \setbeamertemplate{frametitle}{}%
    \input{#1}%
  \end{frame}%
}

% ============================================================
% 正文行与分组
% ============================================================

% --- 正文行：全页只有这一种行 ---
% 姓名、单位、每一条要点，全部用同一个命令，因此字号、行距、
% 字重、颜色必然完全一致 —— 从代码层面就不可能出现不一致。
% 不要为"姓名"或"要点"另建命令：两个命令意味着两套样式，
% 迟早会写出不一样的效果。一个命令，一个样式。
\newcommand{\DeckLine}[1]{%
  \par
  \hangindent=1.1em\hangafter=1
  \noindent #1\par
  \vspace{0.30cm}%
}

% --- 分组留白 ---
% 唯一用来分组的手段：在两组信息之间插入一次更大的空白。
% 不用分隔线、不用小标题、不用字号差异。
\newcommand{\DeckGap}{%
  \par\vspace*{0.40cm}%
}

% --- 带时间标签的行 ---
% #1 时间标签，#2 内容
%
% 不用空格去凑对齐：标签字数不等（"大二寒假" 4 字、"本学期" 3 字），
% 手打空格永远对不齐，而且全角/半角混用会更乱。
% 这里用一个固定宽度的盒子把标签撑到统一宽度，内容自然左对齐。
\newlength{\DeckTagWidth}
\setlength{\DeckTagWidth}{5.6em}   % 约容纳 5 个全角字符
\newcommand{\DeckTagged}[2]{%
  \DeckLine{\makebox[\DeckTagWidth][l]{#1}#2}%
}

% ============================================================
% 条目（带链接）
% ============================================================

% --- 条目：一条两行（标题行 + 说明行）---
% #1 标题行（可内嵌 \DeckLink 让标题本身可点），#2 说明行（角色 / 出处）
% 两行之间比 \DeckLine 更紧（0.12cm），让它们读作一条完整记录；
% 条目之间用 0.34cm 拉开。标题折行时第二行与首行文字左对齐。
\newcommand{\DeckCite}[2]{%
  \par
  \hangindent=1.1em\hangafter=1
  \noindent #1\par
  \vspace{0.12cm}%
  \noindent #2\par
  \vspace{0.34cm}%
}

% --- 超链接 ---
% beamer 已加载 hyperref，直接用 \href。
% 颜色用强调色，不加下划线边框（hyperref 的默认方框很难看）。
\newcommand{\DeckLink}[2]{%
  \href{#1}{\color{SlideAccent}#2}%
}

% ============================================================
% 表格
% ============================================================

% --- 表格行距（可调旋钮）---
% 默认值 = 正文行距，因此表格与正文看起来是一套节奏。
% 某页需要更紧或更松时，在该页文件开头改这一个长度即可：
%     \setlength{\DeckTableLeading}{14pt}
% 注意：字号不在这里调。字号只有两档，全篇一致，不允许逐页改动。
\newlength{\DeckTableLeading}
\setlength{\DeckTableLeading}{\baselineskip}

% --- 数值表：「名称占满 + 数值右对齐」---
% #1 表体：名称 & 数值 \\
% 表宽自动等于 \linewidth。表内文字用 \SlideTextFont，字号与正文一致；
% 层次只靠加粗（列头）与列间距。不画竖线、不画外框。
\newcommand{\DeckCourseTable}[1]{%
  \par\noindent
  {\SlideTextFont
   \setlength{\baselineskip}{\DeckTableLeading}%
   \setlength{\tabcolsep}{0pt}%
   \renewcommand{\arraystretch}{1.0}%
   \noindent
   \begin{tabularx}{\linewidth}{@{}>{\raggedright\arraybackslash}X@{\hspace{0.16cm}}r@{}}
     #1
   \end{tabularx}}%
  \par
}

% --- 列表式两列表：「窄标签左 + 长文字占满左对齐」---
% 与 \DeckCourseTable 的区别只在列宽分配方向：
%   数值表是「名称占满、数值右对齐」；
%   本表是「短标签在左（窄）、长文字在右（占满、左对齐）」。
% 年份 + 描述、时间线这类内容用这个，
% 否则描述会被右对齐得零碎，且长条目会溢出列宽。
\newcommand{\DeckListTable}[1]{%
  \par\noindent
  {\SlideTextFont
   \setlength{\baselineskip}{\DeckTableLeading}%
   \setlength{\tabcolsep}{0pt}%
   \renewcommand{\arraystretch}{1.0}%
   \noindent
   \begin{tabularx}{\linewidth}{@{}l@{\hspace{0.55cm}}>{\raggedright\arraybackslash}X@{}}
     #1
   \end{tabularx}}%
  \par
}

% --- 表头与细线 ---
% 表头行：与正文同字号，只靠加粗区分
\newcommand{\DeckCourseHead}[2]{%
  \textbf{#1} & \textbf{#2} \\[0.04cm]
}
% 列表式表头
\newcommand{\DeckListHead}[2]{%
  \textbf{#1} & \textbf{#2} \\[0.04cm]
}
% 列头与第一行数据之间的细线。
% 用 \noalign 单独发 \hrule，避免 \hline\noalign 给表格带来的异常深度
% （实测那一手会让 5 行表的深度变成 77pt）。
\newcommand{\DeckCourseRule}{%
  \noalign{\vspace{0.09cm}\hrule height 0.4pt\vspace{0.13cm}}%
}

% ============================================================
% 图片
% ============================================================

% --- 两图并排 ---
% #1 左图路径，#2 右图路径，#3 左图高度，#4 右图宽度
%
% 不用 beamer 的 columns：它的栏间距会与给定宽度叠加，很难算准，
% 几次都报 Overfull。改成两图放在同一行、用 \hfill 撑开，
% 宽度完全由 \includegraphics 的 height/width 决定，可控。
%
% \includegraphics 的 width/height 是"缩放后"尺寸：只给 height 时，
% 宽度按原图比例推算，所以两张图若比例不同、都想等高，
% 就必须给其中一张显式 width（调用方算好传入）。
%
% 关键约束：两图宽度之和必须明显小于正文宽（14.25cm），留 0.3cm 以上。
% 只差 0.03cm 都会让第二张图被挤到下一行，并报 overfull vbox。
\newcommand{\DeckFigures}[4]{%
  \par\noindent
  \includegraphics[height=#3]{#1}%
  \hfill
  \includegraphics[width=#4]{#2}%
  \par
}
}{}
